\section{Experimental Results and Empirical Findings}
\label{sec:results}

\subsection{Clean-Data Recognition Baseline}
The proposed Random Forest method obtained clean macro-F1 0.8574 on the configured evaluation split.

\subsection{Fault Resilience Under Simulated Sensor Failure}
For moderate dropout, the proposed method obtained macro-F1 0.7919. Results for other injected faults are tabulated separately.
The \textit{SensorReliabilityEstimator} obtained mean AUROC 0.8882 when clean and injected-fault windows were compared; injected labels were used only for evaluation.

\subsection{Sensing Efficiency and Cost Savings}
The runtime controller responded to reliability, entropy, budget, and periodic probes.
The normalized sensing cost proxy $C(t)$ was 0.53 (a 47.2\% reduction compared to always-active sensing). The all-sensor uniform clean macro-F1 was 0.8891, compared with 0.8574 for the adaptive method.
